%% ============================================================================
%%  Fisheries ecology manuscript draft template
%%  ---------------------------------------------------------------------------
%%  Compile with XeLaTeX or LuaLaTeX (required: fontspec / Libertinus Serif).
%%  Overleaf: Menu > Settings > Compiler = XeLaTeX.
%%  Bibliography: BibTeX + fishfish.bst (place fishfish.bst and bib.bib in
%%  the same folder as this file).
%%
%%  Quick start
%%    1. Fill in the title / author block below.
%%    2. Write in the sections; use \citep{} / \citet{} for references.
%%    3. Toggle \linenumbers on/off for review vs. clean copies.
%%    4. Toggle \doublespacing (in the \begin{spacing} call) as needed.
%% ============================================================================

\documentclass[11pt]{article}

% ---- Encoding & page geometry ----------------------------------------------
\usepackage[utf8]{inputenc}
\usepackage[margin=1.25in]{geometry}

% ---- Math ------------------------------------------------------------------
\usepackage{amsmath, amsfonts, amssymb}
\usepackage{amsbsy}
\usepackage{bm}

% ---- Colour & links --------------------------------------------------------
\usepackage[dvipsnames]{xcolor}
\definecolor{niceblue}{HTML}{236899}
\definecolor{darkgrey}{HTML}{A9A9A9}
% \rev{} = highlight revised text (e.g. in response to reviewers)
\newcommand{\rev}[1]{{\color{niceblue} #1}}
% \ml{} = margin/track comment in red (rename per co-author as you like)
\newcommand{\ml}[1]{{\color{red}#1}}

\usepackage{hyperref}
\usepackage{nameref}
\hypersetup{
    colorlinks=true,
    linkcolor=black,
    filecolor=black,
    urlcolor=niceblue,
    citecolor=niceblue,
    linkbordercolor=white
}

% ---- Figures, floats, tables -----------------------------------------------
\usepackage{epsfig}
\usepackage{subfig}
\usepackage{float}
\usepackage{placeins}     % \FloatBarrier
\usepackage{booktabs}     % \toprule \midrule \bottomrule
\usepackage{pdflscape}    % landscape pages for wide tables/figures
\usepackage{array}

% ---- Text layout helpers ---------------------------------------------------
\usepackage{xspace}
\usepackage{ragged2e}
\usepackage{comment}
% \usepackage{hanging}    % hanging-indent reference list (SI) - not installed
                          % locally and unused elsewhere in this file; re-add
                          % if a hanging-indent SI reference list is needed
\usepackage{setspace}
\usepackage{authblk}      % author / affiliation block

% ---- Line numbers (for review) ---------------------------------------------
\usepackage{lineno}
% Make amsmath align/gather environments play nicely with line numbers:
\makeatletter
\let\LN@align\align
\let\LN@endalign\endalign
\renewcommand{\align}{\linenomath\LN@align}
\renewcommand{\endalign}{\LN@endalign\endlinenomath}
\let\LN@gather\gather
\let\LN@endgather\endgather
\renewcommand{\gather}{\linenomath\LN@gather}
\renewcommand{\endgather}{\LN@endgather\endlinenomath}
\makeatother
\renewcommand\linenumberfont{\normalfont\bfseries\small\color{darkgrey}}
\modulolinenumbers[2]     % show every 2nd line number

% ---- Avoid orphan/widow lines ----------------------------------------------
\widowpenalty10000
\clubpenalty10000

% ---- Bibliography ----------------------------------------------------------
\usepackage[round]{natbib}
\bibliographystyle{fishfish}
\bibpunct{(}{)}{,}{a}{}{,}

% ---- Fonts (requires XeLaTeX/LuaLaTeX) -------------------------------------
\usepackage{fontspec}
\setmainfont{Libertinus Serif}   % or: \setmainfont{Times New Roman}

% ---- Custom column types for nicely wrapped table cells --------------------
\newcolumntype{^}{>{\currentrowstyle}}
\newcommand{\PreserveBackslash}[1]{\let\temp=\\#1\let\\=\temp}
\newcolumntype{C}[1]{>{\PreserveBackslash\centering}p{#1}}
\newcolumntype{R}[1]{>{\PreserveBackslash\raggedleft}p{#1}}
\newcolumntype{L}[1]{>{\PreserveBackslash\raggedright}p{#1}}

% ---- A slightly bolder title font ------------------------------------------
\newcommand*{\TitleFont}{%
      \usefont{\encodingdefault}{\rmdefault}{b}{n}%
      \fontsize{13}{15}%
      \selectfont}

% ---- Cross-reference helpers for reviewer line references ------------------
\newcommand{\R}[1]{\label{#1}\linelabel{#1}}
\newcommand{\lr}[1]{line~\lineref{#1}}

\date{}

% ============================================================================
\begin{document}
\begin{spacing}{2}   % double spacing; change to {1} or use \onehalfspacing

% ---- TITLE / AUTHOR BLOCK --------------------------------------------------
\noindent
\textbf{Journal Name}

\noindent
\textbf{Article Type}   % e.g. Research Article

\noindent
\textbf{Do biodiversity hotspots represent rare fishes? A test across the North Sea–Baltic transition}
%    Species-rich areas do not necessarily represent rare species

\noindent
\textbf{Authors}: Federico Maioli$^{1*}$, P. Daniël van Denderen$^{1}$, Benjamin Weigel$^{2}$, Max Lindmark$^{3}$, Fernanda C. da Silva$^{1}$,  Ole Henriksen$^{1}$, Tobias Mildenberger$^{1}$, Mikael van Deurs$^{1}$, Martin Lindegren$^{1}$
\\~\\
\noindent
$^1$Technical University of Denmark (DTU), National Institute of Aquatic Resources (DTU Aqua), Kgs. Lyngby, Denmark

\noindent
$^2$INRAE, France

\noindent
$^3$Swedish University of Agricultural Sciences (SLU), Sweden
\\~\\
\noindent
\textbf{*Corresponding author}: \href{mailto:name@email.com}{name@email.com}

\noindent
\textbf{ORCID:} First A. Author (0000-0000-0000-0000), Second B. Author (0000-0000-0000-0000).

% Turn line numbers on for review copies; comment out for clean copies.
\linenumbers
\setcounter{secnumdepth}{0}   % remove section numbers
\clearpage

% ---- ABSTRACT --------------------------------------------------------------
\section{Abstract}
% ~150--300 words. State the question, the data/approach, the key findings,
% and why they matter.
Write the abstract here.

\vspace*{\fill}
\noindent
Keywords:
Keyword one,
Keyword two,
Keyword three,
Keyword four,
Keyword five\\
\clearpage

% ---- INTRODUCTION ----------------------------------------------------------
\section{Introduction}
Species richness is one of the most widely used metrics for identifying areas of high biodiversity and informing spatial conservation planning. However, species differ substantially in the size and spatial configuration of their distributions, such that areas supporting many species may disproportionately represent widespread species while failing to capture geographically restricted ones \citep{CoincidenceRarityRichness,ContributionRarityCommonness}. Widespread species exert a disproportionate effect on patterns of species richness relative to range-restricted taxa, so richness hotspots can systematically under-represent the rare and vulnerable species that conservation planning is often intended to protect. Species richness may therefore provide an incomplete, and potentially misleading, proxy for biodiversity value when the planning objective is the representation of rarity rather than of raw species counts.

A parallel but distinct question concerns functional diversity: the range of ecological strategies -- body size, trophic position, habitat use, life history -- present in a community. Richness and functional diversity are not interchangeable: a species-rich assemblage can be functionally redundant if its species share similar ecological roles, while a species-poor assemblage can hold disproportionate functional distinctiveness. Because standard functional diversity indices are computed from the same community data as richness, they are also known to be positively biased by species count for a purely combinatorial reason -- more species sampled means more trait combinations sampled, independent of any ecological signal -- making a null-model correction necessary before functional diversity can be compared against richness on equal terms.

These questions are particularly relevant for marine spatial planning, where environmental gradients generate strong spatial turnover in community composition, and where existing marine protected area (MPA) networks already constrain where any new protection can be added. Approaches incorporating species rarity or functional diversity can identify spatial priorities that differ substantially from those based on richness alone, and the practical consequence of that difference -- which specific sites a reserve of a given size would actually protect -- is best evaluated directly, under a realistic area budget and against the protection already in place, rather than inferred from how similar two metrics look in the aggregate.

Species distribution models fitted to systematic trawl survey data provide a framework for investigating these patterns at scale, by generating spatially standardized predictions of occurrence for many species simultaneously and propagating the uncertainty in those predictions through to derived community metrics \citep{wartonManyVariablesJoint2015, ovaskainenHowMakeMore2017, ovaskainen_abrego_2020}.

Here, we fit species distribution models to three decades of scientific bottom-trawl survey data spanning the continental shelf of the north-east Atlantic, from the Bay of Biscay and Celtic Sea to the Nordic Seas, to ask whether species richness is a reliable proxy for two alternative conservation-relevant properties of a fish community: species rarity and functional diversity. Specifically, we (i) quantify the spatial relationship between species richness, rarity-weighted richness, and functional diversity (both raw and corrected for the richness artifact via a null model), (ii) test how much of each alternative property a richness-based spatial prioritization would actually capture relative to a prioritization targeted at that property directly, under a fixed area budget and with existing MPA coverage held in place, and (iii) identify the specific sites and species a richness-based strategy would systematically miss. By comparing global distributional overlap against budget-constrained prioritization outcomes directly, we also assess whether aggregate similarity between richness and other biodiversity metrics is informative about their value as conservation planning proxies, or whether it obscures decision-relevant disagreement concentrated in the sites that actually compete for protection.
% ---- METHODS ---------------------------------------------------------------
\section{Methods}

\subsection{Study area and survey data}
We compiled bottom-trawl survey data from the ICES Database of Trawl Surveys (DATRAS), covering the continental shelf of the north-east Atlantic from the Bay of Biscay and Celtic Sea, through the Irish Sea, English Channel and North Sea, to the Nordic Seas margin. Two categories of survey were excluded a priori: the Baltic International Trawl Survey (BITS), because the Baltic proper is a distinct, low-salinity system whose species pool and depth associations are not representative of the Atlantic shelf; and two surveys (NS-IDPS, IS-IDPS) that, despite their naming, sample the Barents Sea and Denmark Strait / Greenland waters well outside the continental shelf domain and were too spatially disjunct from the rest of the survey grid to be connected by the spatial model. After these exclusions, the dataset comprised 94,854 hauls from 1993--2025 across 23 surveys. Species identifications were retained at the species level for bony fish, cartilaginous fish, and jawless fish (classes Teleostei, Elasmobranchii, Holocephali, Chondrostei, Petromyzonti and Myxini); records identified only to genus or family were excluded. For each haul, species were scored as present or absent in the catch, and species recorded in fewer than 30 hauls overall were excluded from further analysis as having too little information to support a fitted depth or temperature response, leaving 268 species.

\subsection{Environmental covariates}
Depth was taken from the GEBCO global bathymetric grid, matched to each haul's coordinates and supplemented by the survey's own recorded depth where available. Bottom temperature was obtained from the Copernicus Marine Service global ocean physics reanalysis/forecast product, matched to each haul by month and year, and supplemented by the survey's own recorded bottom temperature where available. All spatial coordinates were projected to the ETRS89-LAEA Europe equal-area projection (EPSG:3035), appropriate given the survey footprint spans from the Nordic Seas to Iberia and a single UTM zone would introduce substantial distortion at that extent.

\subsection{Species distribution models}
We fitted a spatial generalized linear mixed model (binomial family, logit link) separately for each species using \texttt{sdmTMB} \citep{andersonSdmTMB2024}, a stacked-species-distribution-model approach in which each species' occurrence probability is modelled independently rather than jointly. Species were split into two groups by their median presence depth (threshold: 200~m, approximating the shelf break): for shelf species, the linear predictor included second-order polynomials of log-depth and of bottom temperature; for species characteristic of deeper water, temperature was dropped from the formula, since deep water varies little in temperature and the term was not identifiable there. Depth entered the model on the log scale, which an AIC comparison run separately for each species favoured over the raw scale for the majority of species. Raw rather than orthogonal polynomials were used throughout so that the fitted design matrix could be reconstructed identically when predicting onto new data.

Each model included a spatial random field estimated via the SPDE approximation on a triangulated mesh, with a coastline barrier \citep{bakkaBarrierSPDE2019} so that spatial correlation could not propagate across land (for example, between the North Sea and Skagerrak sides of Jutland, which are close in Euclidean distance but separated by the Jutland peninsula). A random intercept by year accounted for interannual variation shared within a species; a fixed year effect was found to be poorly identified for species not recorded in every survey year and was replaced by this random-intercept formulation. Fixed effects were initialized from a maximum-likelihood fit before sampling (\texttt{initPar = "fixed effects"}) to improve convergence.

Model convergence was assessed for every species via gradient-based and Hessian-based diagnostics implemented in \texttt{sdmTMB}'s \texttt{sanity()} check; species failing this check on a first fit were given a short additional optimization pass before being re-evaluated. Of the 268 species fitted, 221 (82\%) converged by this criterion and were retained for all subsequent spatial analyses; the 47 non-converging species were predominantly, but not exclusively, those with very low prevalence, for which the spatial field variance was poorly identified. Predictive performance was evaluated in-sample for every species using the area under the receiver operating characteristic curve (AUC) and the true skill statistic (TSS, sensitivity + specificity $-$ 1) at each species' own optimal classification threshold, since a fixed 0.5 threshold systematically under-credits rare species whose fitted probabilities rarely cross it.

\subsection{Prediction grid}
Predictions were generated on a regular 4$\times$4~km lattice in the same projection as the survey data, retaining only cells within 20~km of an actual haul so that predictions never extrapolated into unsurveyed water. Grid cells with depth or temperature outside the range observed in the training data were also excluded, since the polynomial response terms are not constrained outside their fitted range and would otherwise produce unstable predictions concentrated along the shelf break, where bathymetry changes steeply over short distances.

\subsection{Species richness and its uncertainty}
For each of the 221 species with a converged model, we predicted occurrence probability across the grid. Two complementary representations of species richness uncertainty were produced. First, parameter uncertainty was propagated by simulating 100 draws per species from the joint precision matrix of each fitted model (\texttt{predict(..., nsim = 100)}), summing the simulated probabilities across species within each draw to give a distribution of expected richness per grid cell. Second, observation-level uncertainty was propagated separately by treating each species' fitted probability as a Bernoulli parameter and drawing 100 independent presence/absence realizations per species per cell, summing presences across species within each draw to give a distribution of realized richness counts. The two sources of uncertainty are not interchangeable: the first describes how precisely each species' response was estimated; the second describes the inherent stochasticity of whether an individual species is actually encountered, given that fitted probability, and it was found to be the larger of the two contributions to total uncertainty in richness at any given cell.

\subsection{Functional diversity and a null-model correction for richness}
Species traits (habitat, feeding mode, trophic level, body and fin shape, aspect ratio, offspring size, spawning type, age at maturity, fecundity, growth parameters, and maximum length and age) were compiled from a published trait collection for North Atlantic and north-east Pacific continental shelf fishes \citep{beukhofTraitCollection2019}, matched to species by binomial name. A mixed-type (Gower) dissimilarity among species was computed from these traits and embedded in a Euclidean space via principal coordinates analysis with a Cailliez correction, since Gower dissimilarity is not generally Euclidean.

For each grid cell, each species was weighted by its predicted occurrence probability, normalized to sum to one across species. We computed two community-weighted functional diversity indices from this weighted trait space: Rao's quadratic entropy \citep{raoDiversityDissimilarity1982, bottadukatRaosQuadraticEntropy2005}, the expected trait distance between two individuals drawn at random from the weighted community, and functional dispersion (FDis) \citep{laliberteDistanceBasedFramework2010}, the weighted mean distance of species to the community's weighted centroid in trait space.

Because both indices are known to be positively, and partly mechanically, correlated with species richness -- a community assembled from more species samples more trait combinations largely independent of any ecological signal -- we additionally computed a standardized effect size (SES) for each index against a null model that repeatedly (500 permutations) reassigns species identities to trait values while holding each cell's predicted community composition (and hence its richness) fixed, then expresses the observed index as $(\text{observed} - \text{mean}_{\text{null}}) / \text{sd}_{\text{null}}$. This isolates whether the species actually predicted to co-occur are more or less functionally distinct than a random relabelling at the same richness, independent of richness itself.

\subsection{Rarity-weighted richness}
As a complementary, conservation-oriented alternative to plain richness, each species was assigned a rarity weight equal to the inverse of its total predicted occurrence summed across the grid -- a continuous analogue of inverse range-size weighting \citep{williamsMeasuringBiodiversityComplementary1996}. A rarity-weighted richness value for each cell was then computed as the occurrence-probability-weighted sum of these species weights, so that cells contribute disproportionately to the total when they hold species that are rarely predicted present elsewhere.

\subsection{Comparing richness, rarity, and functional diversity}
We compared the spatial pattern of richness against rarity-weighted richness, raw and SES-corrected functional diversity, and the richness of a subset of species identified as sensitive by the ICES Workshop on Biodiversity Sensitive Species (WKABSENS). Three complementary approaches were used. First, we computed threshold-free overlap between pairs of spatially normalized surfaces using Schoener's $D$ \citep{schoenerNonsynchronousSpatialOverlap1970} and Bhattacharyya's coefficient \citep{bhattacharyyaMeasureDivergence1943}. Second, we evaluated budget-constrained performance directly: ranking grid cells by a given metric, we asked what share of the total rarity-weighted (or functional-diversity) value a reserve of a given size would capture, relative to a random selection and to the best possible (optimal) selection ranked by that value itself, across reserve sizes from 1\% to 50\% of the study area. Third, we evaluated species-level representation: for a reserve of a given size, how many species have at least 50\% of their predicted total occurrence inside it, under richness-based versus rarity-based selection, identifying species left essentially unprotected under one strategy despite being well represented under the other.

\subsection{Spatial prioritization against existing protection}
We obtained marine protected area boundaries from the World Database on Protected Areas \citep{unepwcmcProtectedPlanet2026}, retaining designated (not proposed) sites with a non-zero marine component and intersecting them with the grid to flag cells already under protection. For a fixed area budget of 30\% of the study area, we constructed a prioritization solution for each candidate metric (richness, rarity-weighted richness, raw and SES-corrected Rao's Q and FDis, and WKABSENS-species richness) by locking in all already-protected cells and filling the remaining budget with the highest-ranked unprotected cells under that metric; with a single budget constraint and no connectivity or cost term, this greedy fill is equivalent to the optimal solution. Agreement between the resulting solutions was quantified with the Jaccard index, computed on the discretionary (non-locked-in) cells only, since every solution trivially agrees on the cells already protected regardless of the ranking metric used.

\subsection{Software}
Analyses were conducted in \textsf{R} \citep{rcoreteamLanguageEnvironment2026}. Species distribution models were fitted with \texttt{sdmTMB} \citep{andersonSdmTMB2024}; spatial data processing used \texttt{sf}; the trait dissimilarity and ordination used \texttt{cluster}; network visualization used \texttt{tidygraph} and \texttt{ggraph}.

% ---- RESULTS ---------------------------------------------------------------
\section{Results}
% draft based on the completed analysis pipeline; figure/table numbers are
% placeholders until the corresponding plots are copied into manuscript/figures/

\subsection{Model performance}
Of the 268 species fitted, 221 (82\%) met the convergence criteria and were retained for all spatial analyses. Discrimination was high across the retained species: median AUC was 0.990 and median TSS at each species' own optimal threshold was 0.925, with no strong relationship between a species' predictive performance and whether its model had converged -- the majority of non-converging species still had well-behaved point predictions, with the unidentified parameter specifically confined to the spatial variance term rather than the fixed effects.

\subsection{Spatial pattern of richness, rarity, and functional diversity}
Predicted species richness averaged 15.5 species per grid cell (range 2.2--30.3) and peaked consistently along the continental shelf break, from the Celtic Sea and Bay of Biscay margin to west of Ireland, with lower richness through the central and eastern North Sea (Fig.~\ref{fig:richness_maps}). Richness computed from a subset of 163 species with available trait data tracked the full-community pattern closely (Pearson's $r = 0.94$); richness of the 59--61 species flagged as conservation-sensitive by WKABSENS tracked it more loosely ($r = 0.82$) and showed visibly lower spatial contrast, consistent with sensitive species being a distinct rather than a scaled-down subset of the full community.

Raw Rao's Q and FDis were both positively correlated with richness ($r = 0.32$ and $0.45$, respectively), and standardizing both indices against the trait-shuffling null model reduced these correlations to approximately zero ($r = -0.04$ and $-0.05$), confirming that most of the apparent spatial covariation between functional diversity and richness is attributable to the combinatorial richness artifact rather than an independent ecological signal. Standardized effect sizes were negative across nearly the entire study area (median SES $\approx -3.1$ for both indices, with positive values in under 0.1\% of cells), indicating that the species predicted to co-occur are consistently more functionally similar to one another than a random reassignment of traits at the same richness would produce -- a signature of environmental filtering rather than functional niche partitioning, present essentially everywhere in the domain rather than concentrated in particular sites.

\subsection{Global overlap versus budget-constrained agreement}
Despite these differences, richness, rarity-weighted richness, and raw functional diversity all showed high global distributional overlap by Schoener's $D$ (range 0.73--0.98 across all pairs) and Bhattacharyya's coefficient (0.95--1.00). Decomposing this overlap showed why: cells in the top 3.6\% of either richness or rarity-weighted value -- the slice that actually competes for a constrained conservation budget -- contributed only 24\% of the total disagreement signal the overlap statistics are built from, with the remaining 76\% coming from the broad middle of the distribution where metrics broadly track one another. Aggregate similarity between richness and other biodiversity metrics is therefore not informative about their relative value as spatial planning proxies; the decision-relevant comparison requires evaluating agreement specifically among the sites that would actually be selected.

\subsection{What richness-based protection would capture}
At a reserve covering the top 10\% of cells by rank, richness-based selection captured 13.2\% of total rarity-weighted value, an efficiency of 43\% relative to the best possible (rarity-ranked) selection of the same size; efficiency rose to 71\% at a 30\% reserve. Against raw functional diversity, richness performed far better (89--99\% efficiency across reserve sizes), but this advantage was again attributable to the richness--functional-diversity confound: against the null-model-corrected indices, richness-based selection achieved only 81--82\% efficiency, roughly midway between its performance for rarity and its (misleadingly high) performance for the uncorrected functional diversity indices.

At the species level, a 10\% richness-based reserve left every one of the 221 species with less than half of its predicted range protected; a rarity-based reserve of the same size adequately protected (defined as $\geq$50\% of predicted range) 59 species. This gap persisted at larger reserve sizes (20\%: 16 versus 73 species; 30\%: 53 versus 103 species). At a 20\% reserve, four species -- \textit{Atherina presbyter}, \textit{Epigonus denticulatus}, \textit{Lepidopus caudatus}, and \textit{Malacocephalus laevis} -- had under 10\% of their predicted range protected under richness-based selection despite being well represented ($\geq$50\%) under rarity-based selection.

\subsection{Prioritization against existing protected areas}
Existing marine protected areas, filtered to designated sites with a marine component, already covered 26.4\% of the study area, close to the 30\% target budget used in the prioritization exercise and leaving a discretionary margin of only 3.6\% of the total area for any metric to add. Within that discretionary margin, richness showed almost no agreement with any other candidate metric (Jaccard index 0.001--0.008 against rarity-weighted richness and raw/SES-corrected Rao's Q and FDis, 0.14 against WKABSENS-species richness), while rarity-weighted richness, raw and SES-corrected functional diversity, and WKABSENS-species richness agreed with one another at substantially higher levels (Jaccard 0.25--0.97). Of the cells added to the reserve under at least one metric, only nine were chosen by all seven metrics simultaneously (Fig.~\ref{fig:prioritization_network}).

% ---- DISCUSSION ------------------------------------------------------------
\section{Discussion}
Interpret the results, relate them to the literature, note limitations, and
close with the broader implications.

% ---- BACK MATTER -----------------------------------------------------------
\section{Acknowledgements}
Thank contributors, funders, and data providers here.

\section{Author contributions}
Describe contributions (e.g. CRediT roles).

\section{Conflicts of interest}
The authors declare that they have no conflict of interest.

\section{Data Availability Statement}
Code and data are available at \url{https://github.com/user/repo} and archived
at \url{https://doi.org/xxxxx}.

\section{Funding}
Funding sources and grant numbers.

\bibliography{bib}

\end{spacing}
\clearpage

% ============================================================================
%  TABLES
% ============================================================================
\section{Tables}

\begin{table}[htp]
  \centering
  \caption{Example table caption.}
  \renewcommand\arraystretch{2}
  \label{tab:example}
  \begin{tabular}{L{3cm} L{3cm} L{3cm}}
    \toprule
    \textbf{Column A} & \textbf{Column B} & \textbf{Column C} \\
    \midrule
    Row 1 & value & value \\
    Row 2 & value & value \\
    \bottomrule
  \end{tabular}
\end{table}
\clearpage

% ============================================================================
%  FIGURES
% ============================================================================
\section{Figures}

\begin{figure}[htb]
    \centering
    \includegraphics[width=0.95\textwidth]{figures/example.pdf}
    \caption{Example figure caption.}
    \label{fig:example}
\end{figure}
\clearpage

% ============================================================================
%  SUPPORTING INFORMATION
%  Figures/tables are relabelled S1, S2, ... and page/line counters reset.
% ============================================================================
\renewcommand{\thetable}{S\arabic{table}}
\setcounter{table}{0}
\renewcommand{\thefigure}{S\arabic{figure}}
\setcounter{figure}{0}
\addtocontents{toc}{\protect\setcounter{tocdepth}{0}}
\renewcommand{\theHfigure}{S\arabic{figure}}
\renewcommand{\theHtable}{S\arabic{table}}
\onehalfspacing
\resetlinenumber
\setcounter{page}{1}
\setcounter{equation}{0}
\nolinenumbers

\begin{Center}
\section{Supporting Information}
\end{Center}

\begin{figure}[htb]
    \centering
    \includegraphics[width=0.65\textwidth]{figures/supp/example_supp.pdf}
    \caption{Example supporting figure.}
    \label{fig:example_supp}
\end{figure}
\FloatBarrier
\clearpage

\end{document}
